\documentclass[12pt]{article}
\usepackage[utf8]{inputenc}
\usepackage[T1]{fontenc}
\usepackage[margin=1in]{geometry}
\usepackage{setspace}
\usepackage{xcolor}
\usepackage[colorlinks=true, urlcolor=blue!60!black, linkcolor=blue!60!black]{hyperref}
\setlength{\parindent}{2em}
\setlength{\parskip}{1.25em}
\renewcommand{\baselinestretch}{1.0}
\begin{document}
\begin{center}
{\LARGE \textit{McCulloch v. Maryland}, 17 U.S. (4 Wheat.) 316 (1819)\par}
\vspace{0.8em}
{\large Case Brief\par}
\end{center}
\vspace{1em}

\section{Memory Jogger}

In \textit{\href{https://www.courtlistener.com/opinion/85272/mculloch-v-state-of-maryland/}{McCulloch v. Maryland}}, the Supreme Court held that the Necessary \& Proper Clause grants Congress implied authority to employ any appropriate, unprohibited means---including chartering a national bank---to execute its enumerated powers, and that the Supremacy Clause forbids states from taxing or impeding federal instrumentalities.

\section{Detailed Case Facts}

On April 10, 1816, following the fiscal disruptions of the War of 1812 after the First Bank of the United States (chartered in 1791) had expired, Congress enacted ``an act to incorporate the subscribers to the Bank of the United States,'' chartering the Second Bank of the United States (\href{https://supreme.justia.com/cases/federal/us/17/316/}{17 U.S. at 317}). The Bank organized in Philadelphia and, in 1817, opened a branch---an ``office of discount and deposit''---in Baltimore, Maryland, without seeking or receiving any charter or authorization from the State of Maryland (\href{https://supreme.justia.com/cases/federal/us/17/316/}{17 U.S. at 318}).

Resentful of the national bank's competition with state-chartered banks and its contraction of credit during the economic downturn of 1818, the Maryland General Assembly passed a statute on February 11, 1818, titled ``An act to impose a tax on all banks or branches thereof, in the State of Maryland not chartered by the legislature'' (\href{https://supreme.justia.com/cases/federal/us/17/316/}{17 U.S. at 317, 320}; \href{https://www.encyclopedia.com/law/legal-and-political-magazines/mcculloch-v-maryland-1819}{Encyclopedia.com}). Although phrased generally, the Baltimore branch of the Second Bank of the United States was the only bank in Maryland not chartered by the state legislature (\href{https://www.annenbergclassroom.org/resource/the-pursuit-of-justice/pursuit-justice-chapter-2-national-bank-federalism/}{Annenberg Classroom}). Effective May 1, 1818, the statute prohibited any unchartered bank branch from issuing notes except on stamped paper furnished by the Treasurer of the Western Shore at steep per-note fees (e.g., 10 cents on a \$5 note, \$1 on a \$100 note, and \$20 on a \$1,000 note), unless the institution relieved itself by paying an annual lump sum of \$15,000 in advance (\href{https://supreme.justia.com/cases/federal/us/17/316/}{17 U.S. at 320--21}). Violating officers forfeited \$500 per offense, recoverable in an action of debt with half going to the informer and half to the State (\href{https://supreme.justia.com/cases/federal/us/17/316/}{17 U.S. at 321--22}).

James William McCulloch, cashier of the Baltimore branch, issued unstamped bank notes to Baltimore resident George Williams in part payment of a discounted promissory note without paying the \$15,000 annual fee (\href{https://supreme.justia.com/cases/federal/us/17/316/}{17 U.S. at 318--19}). John James, suing as a \textit{qui tam} informer on behalf of himself and the State of Maryland, brought an action of debt against McCulloch to recover the statutory penalties (\href{https://supreme.justia.com/cases/federal/us/17/316/}{17 U.S. at 317}).

\section{Procedural History}

\begin{enumerate}
    \item \textbf{Baltimore County Court (Maryland, 1818):} The parties submitted an agreed statement of facts stipulating that if the Maryland tax statute were valid, judgment should be entered for John James and the State for \$2,500 plus costs; if invalid, judgment of \textit{non pros} should be entered for McCulloch (\href{https://supreme.justia.com/cases/federal/us/17/316/}{17 U.S. at 319--20}). The Baltimore County Court entered judgment against McCulloch (\href{https://supreme.justia.com/cases/federal/us/17/316/}{17 U.S. at 317}).
    \item \textbf{Maryland Court of Appeals (1818):} Maryland's highest court affirmed the judgment of the Baltimore County Court, upholding the state tax (\href{https://supreme.justia.com/cases/federal/us/17/316/}{17 U.S. at 317}; \href{https://www.oyez.org/cases/1789-1850/17us316}{Oyez}).
    \item \textbf{U.S. Supreme Court (1819):} McCulloch sought review by writ of error under Section~25 of the Judiciary Act of 1789. After nine days of oral argument (February 22 to March 3, 1819) featuring leading advocates including Daniel Webster, William Pinkney, and Attorney General William Wirt for McCulloch, and Luther Martin, Joseph Hopkinson, and Walter Jones for Maryland, the Supreme Court issued its unanimous decision three days later on March 6, 1819, reversing the Maryland Court of Appeals and the Baltimore County Court and ordering judgment entered for McCulloch (\href{https://supreme.justia.com/cases/federal/us/17/316/}{17 U.S. at 316, 437}).
\end{enumerate}

\section{Judicial Votes}

Chief Justice John Marshall delivered the unanimous opinion of the Court, writing, ``We are unanimously of opinion that the law passed by the Legislature of Maryland, imposing a tax on the Bank of the United States is unconstitutional and void'' (\href{https://supreme.justia.com/cases/federal/us/17/316/}{17 U.S. at 436}):

\begin{itemize}
    \item \textbf{Unanimous Opinion of the Court (7--0 in standard databases; 6--0 per historical attendance records):} Chief Justice John Marshall (author), joined by Associate Justices Bushrod Washington, William Johnson, Henry Brockholst Livingston, Gabriel Duvall, and Joseph Story. Standard case databases including \href{https://www.oyez.org/cases/1789-1850/17us316}{Oyez}, \href{https://supreme.justia.com/cases/federal/us/17/316/}{Justia}, and \href{https://www.annenbergclassroom.org/resource/the-pursuit-of-justice/pursuit-justice-chapter-2-national-bank-federalism/}{Annenberg Classroom} list Associate Justice Thomas Todd as joining a 7--0 decision, whereas several historical reference works (e.g., \href{https://www.encyclopedia.com/law/legal-and-political-magazines/mcculloch-v-maryland-1819}{Encyclopedia.com}) note that Justice Todd was absent due to illness and did not participate, making the participating vote 6--0.
    \item \textbf{Concurring:} None.
    \item \textbf{Dissenting:} None.
\end{itemize}

\section{Holding}

\textbf{Judgment:} Reversed and annulled the judgments of both the Maryland Court of Appeals and the Baltimore County Court, and ordered judgment entered in the Baltimore County Court for James W. McCulloch (\href{https://supreme.justia.com/cases/federal/us/17/316/}{17 U.S. at 437}).

\textbf{Rule Announced:}
\begin{enumerate}
    \item \textbf{Implied Powers under the Necessary \& Proper Clause (Art.\ I, \S~8, cl.\ 18):} Congress possesses implied powers to choose any appropriate, unprohibited means plainly adapted to executing its enumerated constitutional powers. ``Necessary'' does not mean strictly indispensable; it encompasses means that are convenient, useful, or conducive to carrying out enumerated ends, and the degree of necessity is a question of legislative discretion rather than judicial review---so long as Congress does not use its powers as a ``pretext'' to accomplish objects not entrusted to the national government (\href{https://supreme.justia.com/cases/federal/us/17/316/}{17 U.S. at 413--23}). Under this standard, Congress had authority to charter the Second Bank of the United States and authorize its branches.
    \item \textbf{Intergovernmental Tax Immunity under the Supremacy Clause (Art.\ VI, cl.\ 2):} States have no power, by taxation or otherwise, to retard, impede, burden, or in any manner control the operations of constitutional federal laws or instrumentalities. Accordingly, Maryland's tax on the operations of the Baltimore branch of the Bank of the United States is unconstitutional and void (\href{https://supreme.justia.com/cases/federal/us/17/316/}{17 U.S. at 436--37}).
\end{enumerate}

\textbf{Holding vs.\ Dicta:}
\begin{itemize}
    \item \textbf{Binding Holding:} (1) The Act of April 10, 1816 incorporating the Second Bank of the United States, including its authority to establish branches in the states without state consent, is a valid exercise of Congress's enumerated fiscal and war powers combined with the Necessary \& Proper Clause; and (2) Maryland's discriminatory tax on the Bank's note-issuing operations violates the Supremacy Clause (\href{https://supreme.justia.com/cases/federal/us/17/316/}{17 U.S. at 424--25, 436--37}).
    \item \textbf{Dicta:} Chief Justice Marshall's famous ``pretext'' limitation---stating that if Congress ``under the pretext of executing its powers, pass laws for the accomplishment of objects not intrusted to the Government,'' the Court would be duty-bound to invalidate them (\href{https://supreme.justia.com/cases/federal/us/17/316/}{17 U.S. at 423})---was not necessary to uphold the Bank. Likewise, Marshall's closing carve-out clarifying that the Court's ruling does not bar nondiscriminatory state taxes on the Bank's real property or on Maryland citizens' proprietary shares in the Bank (\href{https://supreme.justia.com/cases/federal/us/17/316/}{17 U.S. at 436}), as well as his suggestion that Congress could constitutionally tax state-chartered banks because all states are represented in Congress (\href{https://supreme.justia.com/cases/federal/us/17/316/}{17 U.S. at 435--36}), were clarifying dicta that nonetheless proved highly influential in later intergovernmental-immunity doctrine.
\end{itemize}

\section{Analysis of Opinions}

\subsection{Unanimous Opinion of the Court (Chief Justice Marshall)}

Chief Justice Marshall divided the opinion into two great constitutional questions: (1) whether Congress had power to incorporate a bank, and (2) whether Maryland could constitutionally tax its Baltimore branch.

\textbf{Part I---Congress's Power to Charter the Bank:} Marshall began with constitutional liquidation through historical practice: the First Congress had intensely debated and chartered the First Bank in 1791, Presidents and successive legislatures had recognized it, and even former skeptics (implicitly President James Madison) had been convinced of its necessity by the fiscal crises following its expiration (\href{https://supreme.justia.com/cases/federal/us/17/316/}{17 U.S. at 401--02}). Turning to first principles, Marshall rejected Maryland's compact theory: the Constitution was ratified not by state legislatures, but by conventions of the people assembled in each state, making the federal government ``emphatically and truly, a Government of the people'' that is ``supreme within its sphere of action'' (\href{https://supreme.justia.com/cases/federal/us/17/316/}{17 U.S. at 403--06}).

Although ``bank'' and ``incorporation'' do not appear among Congress's enumerated powers, Marshall stressed two textual clues. First, unlike Article~II of the Articles of Confederation, the Tenth Amendment deliberately omits the word ``expressly'' before ``delegated,'' leaving room for incidental powers (\href{https://supreme.justia.com/cases/federal/us/17/316/}{17 U.S. at 406--07}). Second, a constitution cannot contain ``the prolixity of a legal code''; it marks ``great outlines,'' leaving the minor means to be deduced from the great powers of ``[t]he sword and the purse''---taxing, borrowing money, regulating commerce, and conducting war across a vast continent (\href{https://supreme.justia.com/cases/federal/us/17/316/}{17 U.S. at 407--08}). Creating a corporation is not a ``great substantive and independent power,'' but merely an ordinary means to execute other sovereign ends (\href{https://supreme.justia.com/cases/federal/us/17/316/}{17 U.S. at 411}).

Marshall then dismantled Maryland's narrow reading of the Necessary \& Proper Clause through four interlocking arguments:
\begin{enumerate}
    \item \textbf{Ordinary Meaning and Intratextual Contrast:} In common usage, ``necessary'' often means ``convenient, or useful, or essential,'' not absolutely indispensable. Tellingly, when the Framers meant strict indispensability in Article~I, Section~10 (limiting state import/export duties to what is ``\textit{absolutely} necessary'' for inspection laws), they added the qualifier ``absolutely''---a modifier omitted from Article~I, Section~8, Clause~18 (\href{https://supreme.justia.com/cases/federal/us/17/316/}{17 U.S. at 413--15}; see \href{https://akhilamar.com/wp-content/uploads/2022/04/Intratextualism.pdf}{Akhil Reed Amar, \textit{Intratextualism}}).
    \item \textbf{Practical Workability:} A strict-indispensability rule would cripple the federal government, invalidating even federal criminal punishments for mail robbery or perjury in federal courts, since post offices and courts could technically exist without criminal penalties (\href{https://supreme.justia.com/cases/federal/us/17/316/}{17 U.S. at 416--18}).
    \item \textbf{Placement and Purpose:} The Necessary \& Proper Clause sits among the \textit{grants} of congressional power in Section~8, not among the \textit{limitations} in Section~9, and its phrasing purports to enlarge or remove doubt about implied powers rather than to abridge them (\href{https://supreme.justia.com/cases/federal/us/17/316/}{17 U.S. at 419--21}).
    \item \textbf{Deference to Legislative Discretion:} Synthesizing his famous test---``Let the end be legitimate, let it be within the scope of the Constitution, and all means which are appropriate, which are plainly adapted to that end, which are not prohibited, but consist with the letter and spirit of the Constitution, are Constitutional''---Marshall concluded that once a measure is appropriate and not pretextual, the degree of its necessity is for Congress, not the courts, to judge (\href{https://supreme.justia.com/cases/federal/us/17/316/}{17 U.S. at 421--23}).
\end{enumerate}

\textbf{Part II---Invalidity of Maryland's Tax:} On the second issue, Marshall reasoned from structural supremacy and political-representation theory. While states retain a broad concurrent power of taxation, that power extends only to subjects within state sovereignty, not to national instrumentalities created by the people of all the states (\href{https://supreme.justia.com/cases/federal/us/17/316/}{17 U.S. at 425--30}). Because ``the power to tax involves the power to destroy,'' permitting a single state to tax a federal instrument would allow a part to control the whole, without the political check of representation: Maryland legislators are accountable only to Maryland voters, not to the entire nation that created and relies on the Bank (\href{https://supreme.justia.com/cases/federal/us/17/316/}{17 U.S. at 428--32, 435--36}). Refusing to entangle courts in line-drawing over how heavy a state tax is ``too much,'' Marshall announced a bright-line bar against state taxation of federal operations, while noting that nondiscriminatory taxes on the Bank's real property or on Maryland citizens' shares remained permissible (\href{https://supreme.justia.com/cases/federal/us/17/316/}{17 U.S. at 430, 436--37}).

\subsection{Concurrences}

There were no concurring opinions.

\subsection{Dissents}

There were no dissenting opinions.

\section{Examples: Future Applications}

\begin{enumerate}
    \item \textbf{Same-Side Application (Federal Civil Commitment of Sexually Dangerous Federal Prisoners):} Suppose Congress authorizes federal courts to civilly commit mentally ill, sexually dangerous federal prisoners beyond the expiration of their federal criminal sentences when states decline to assume custody. Under \textit{McCulloch}, even though the Constitution does not expressly mention prisons or civil commitment, Congress has implied power to criminalize conduct tied to enumerated powers, operate federal prisons, and adopt reasonably adapted custodial measures to protect communities from dangerous prisoners in federal care---the result reached in \textit{\href{https://supreme.justia.com/cases/federal/us/560/126/}{United States v. Comstock}}, 560 U.S. 126 (2010) (and extended to former military offenders in \textit{\href{https://supreme.justia.com/cases/federal/us/570/387/}{United States v. Kebodeaux}}, 570 U.S. 387 (2013)).

    \item \textbf{Same-Side Application (Discriminatory State Tax Targeting Federal Operations):} Suppose a state enacts a special ``federal operations surcharge'' of \$50 per parcel on deliveries made by the U.S. Postal Service within the state, while exempting private couriers. Under \textit{McCulloch}'s second holding, the tax is flatly unconstitutional under the Supremacy Clause because it directly taxes and burdens the operations of a federal instrumentality created pursuant to Congress's enumerated postal power, without political representation of the national citizenry in the state legislature.

    \item \textbf{Opposite-Side Application (Compelling Individuals to Enter Commerce to Facilitate a Regulatory Scheme):} Suppose Congress enacts a mandate requiring uninsured individuals to purchase private health insurance, defending it as ``necessary and proper'' to make guaranteed-issue insurance regulations under the Commerce Clause work smoothly. Under Chief Justice Roberts's controlling analysis in \textit{\href{https://supreme.justia.com/cases/federal/us/567/11-393/case.pdf}{NFIB v. Sebelius}}, 567 U.S. 519 (2012), while the mandate may be ``necessary'' (useful) to the regulatory scheme, it is not ``proper'' under \textit{McCulloch} because it exercises a ``great substantive and independent power''---compelling individuals into commercial activity---rather than an incidental means consonant with the Constitution's structure of limited, enumerated powers.

    \item \textbf{Opposite-Side Application (Nondiscriminatory State Property Tax on Federal Bank Real Estate or Employee Income):} Suppose a state levies a general, nondiscriminatory real-estate tax on all commercial office buildings in the state---including a building owned by a federally chartered bank---or levies a general state income tax on the salaries of federal employees residing in the state. Under \textit{McCulloch}'s explicit carve-out (\href{https://supreme.justia.com/cases/federal/us/17/316/}{17 U.S. at 436}), such general taxes are valid because they fall evenly on the state's own constituents (providing the political check Marshall emphasized) and do not target or impede the operations of the federal government.

    \item \textbf{Fence-Sitter (Multi-Link Chain of Implied Powers Reaching Traditional State Police Powers):} Suppose Congress invokes the Necessary \& Proper Clause in combination with the Spending Clause to enact a federal criminal statute making it a felony for any private employee of a local contractor---who receives indirect subgrants from a state agency that accepts federal highway funds---to falsify any local business record, even where no federal funds are directly stolen or misapplied. Under \textit{McCulloch}, this is a genuine fence-sitter: proponents would argue it is a convenient, rational means to safeguard the integrity of federally funded programs (analogous to \textit{Comstock} and \textit{Sabri v. United States}), while challengers would contend that stacking multiple inferences stretches ``plainly adapted'' too far, trenches on traditional state criminal jurisdiction, and crosses Marshall's line against exercising an unenumerated ``great substantive and independent power'' or acting under a ``pretext.''
\end{enumerate}

\section{Critique}

\textbf{Originalist and Textualist Perspectives:} Immediately after \textit{McCulloch} was decided, Virginia Judge Spencer Roane (writing as ``Amphictyon'' and ``Hampden'') attacked Marshall's opinion for dissolving enumerated limits on Congress and arrogating to the federal judiciary exclusive authority to police federalism, prompting Marshall's extraordinary pseudonymous defense as ``A Friend of the Constitution'' (\href{https://www.cambridge.org/core/books/monitoring-american-federalism/state-challenges-to-the-supreme-courts-control-over-constitutional-interpretation/85539C2F1811467C31EF62DBD350C7D8}{Cambridge}; \href{https://www.annenbergclassroom.org/resource/the-pursuit-of-justice/pursuit-justice-chapter-2-national-bank-federalism/}{Annenberg Classroom}). Modern originalist scholars have sharpened this critique. Drawing on eighteenth-century agency, fiduciary, and corporate law, \href{https://scholarship.law.bu.edu/cgi/viewcontent.cgi?filename=0&article=4259&context=faculty_scholarship&type=additional}{Gary Lawson, Geoffrey Miller, Robert Natelson, and Guy Seidman} argue that Marshall erred by equating ``necessary'' with merely ``convenient'' or ``useful'' and by collapsing ``necessary and proper'' into deferential rational-basis review; in their view, ``necessary'' required a tighter telic/incidental fit under agency law, while ``proper'' imposed independent fiduciary and jurisdictional limits (such as respecting federalism and separation of powers). Similarly, \href{https://harvardlawreview.org/forum/vol-128/sharing-the-necessary-and-proper-clause/}{William Baude} highlights James Madison's distinction between ``incidental'' powers and ``great substantive and independent powers,'' a boundary revived in \textit{\href{https://supreme.justia.com/cases/federal/us/567/11-393/case.pdf}{NFIB v. Sebelius}}. Conversely, originalist and intratextualist defenders such as \href{https://akhilamar.com/wp-content/uploads/2022/04/Intratextualism.pdf}{Akhil Reed Amar} and \href{https://constitutioncenter.org/the-constitution/articles/article-i/clauses/754}{Jack Balkin} praise Marshall's intratextual contrast with ``absolutely necessary'' in Article~I, Section~10, and argue that ``necessary'' originally encompassed convenient means for organizing and executing federal powers.

\textbf{Progressive, Doctrinal, and Structural Perspectives:} Progressive and nationalist scholars---including \href{https://scholarship.law.umn.edu/cgi/viewcontent.cgi?article=2232&context=concomm}{John Mikhail} and \href{https://www.law.georgetown.edu/public-policy-journal/wp-content/uploads/sites/23/2021/09/Schwartz.pdf}{David Schwartz}---argue that \textit{McCulloch} is if anything \textit{too} deferential to strict ``enumerationism.'' Mikhail emphasizes that the Necessary \& Proper Clause's ``all other Powers vested by this Constitution in the Government of the United States'' language (the ``Sweeping Clause'') points back to corporate and national powers vested in the United States as a whole, including those reflected in the Preamble, while Schwartz contends that Marshall's concession that the federal government is strictly one of ``enumerated powers'' sat in uneasy tension with his own broad functional nationalism. Doctrinally, \textit{McCulloch} became the bedrock for the modern administrative state, New Deal economic regulation, and federal criminal law (\href{https://www.annenbergclassroom.org/resource/the-pursuit-of-justice/pursuit-justice-chapter-2-national-bank-federalism/}{Annenberg Classroom}), while its political-process rationale for intergovernmental tax immunity (barring taxation without representation of the national public) remains one of the most influential structural arguments in constitutional law (\href{https://akhilamar.com/wp-content/uploads/2022/04/Intratextualism.pdf}{Amar}).

\textbf{Competing Constitutional Values:} \textit{McCulloch} balances \textit{effective, adaptable national governance} (empowering a continental republic to meet unforeseen crises without constant constitutional amendment) against \textit{limited, enumerated federal authority and state fiscal autonomy} (preventing implied powers from swallowing the Tenth Amendment's reservation of powers to the states and the people).

\section{Key Quotations}

\begin{quote}
``Its nature, therefore, requires that only its great outlines should be marked, its important objects designated, and the minor ingredients which compose those objects be deduced from the nature of the objects themselves. \ldots{} In considering this question, then, we must never forget that it is a Constitution we are expounding.''
\end{quote}
\noindent --- Chief Justice John Marshall, Opinion of the Court, \href{https://supreme.justia.com/cases/federal/us/17/316/}{17 U.S. (4 Wheat.) at 407} (also at \href{https://www.courtlistener.com/opinion/85272/mculloch-v-state-of-maryland/}{CourtListener}).

\begin{quote}
``This provision is made in a Constitution intended to endure for ages to come, and consequently to be adapted to the various crises of human affairs.''
\end{quote}
\noindent --- Chief Justice John Marshall, Opinion of the Court, \href{https://supreme.justia.com/cases/federal/us/17/316/}{17 U.S. (4 Wheat.) at 415} (also at \href{https://www.courtlistener.com/opinion/85272/mculloch-v-state-of-maryland/}{CourtListener}).

\begin{quote}
``Let the end be legitimate, let it be within the scope of the Constitution, and all means which are appropriate, which are plainly adapted to that end, which are not prohibited, but consist with the letter and spirit of the Constitution, are Constitutional.''
\end{quote}
\noindent --- Chief Justice John Marshall, Opinion of the Court, \href{https://supreme.justia.com/cases/federal/us/17/316/}{17 U.S. (4 Wheat.) at 421} (also at \href{https://www.courtlistener.com/opinion/85272/mculloch-v-state-of-maryland/}{CourtListener}).

\begin{quote}
``That the power to tax involves the power to destroy; that the power to destroy may defeat and render useless the power to create; that there is a plain repugnance in conferring on one Government a power to control the constitutional measures of another, which other, with respect to those very measures, is declared to be supreme over that which exerts the control, are propositions not to be denied.''
\end{quote}
\noindent --- Chief Justice John Marshall, Opinion of the Court, \href{https://supreme.justia.com/cases/federal/us/17/316/}{17 U.S. (4 Wheat.) at 431} (also at \href{https://www.courtlistener.com/opinion/85272/mculloch-v-state-of-maryland/}{CourtListener}).

\begin{quote}
``The result is a conviction that the States have no power, by taxation or otherwise, to retard, impede, burden, or in any manner control the operations of the constitutional laws enacted by Congress to carry into execution the powers vested in the General Government.''
\end{quote}
\noindent --- Chief Justice John Marshall, Opinion of the Court, \href{https://supreme.justia.com/cases/federal/us/17/316/}{17 U.S. (4 Wheat.) at 436} (also at \href{https://www.courtlistener.com/opinion/85272/mculloch-v-state-of-maryland/}{CourtListener}).

\end{document}
